\chapter{{\fontsize{14pt}{21pt}\selectfont\MakeUppercase{Техническая реализация системы}}}
\label{cha:design}

\section{Архитектура системы}

Разрабатываемая система реализована как модульный монолит с REST API, объединённый в одном репозитории, но логически разделённый на серверную часть, веб-клиент и мобильный клиент. Такой подход упрощает сопровождение учебного проекта, позволяет явно показать взаимодействие подсистем и при этом сохраняет понятную декомпозицию по функциональным модулям.\cite{martin_clean_architecture_ru,fowler_patterns_ru,sommerville2019software}

Серверная часть выступает единственным источником бизнес-логики и данных. Она отвечает за аутентификацию, выбор текущей организации, проверку ролей, работу с курсами и уроками, управление тестами, выполнение адаптивного тестирования, формирование рекомендаций и подготовку аналитики. Веб-клиент ориентирован на роли администратора и преподавателя. Мобильный клиент ориентирован на роль слушателя и закрывает сценарии прохождения обучения.

На рисунке~\ref{fig:design-context} представлена контейнерная схема системы. Диаграмма показывает состав программного комплекса на уровне контейнеров и основные каналы взаимодействия между web-клиентом, mobile-клиентом, REST API серверной части, базой данных и хранилищем медиафайлов.\cite{booch2006uml}

\begin{figure}[H]
\centering
\includegraphics[width=\textwidth,height=0.8\textheight,keepaspectratio]{c4}
\caption{Контейнерная схема программного комплекса}
\label{fig:design-context}
\end{figure}

С точки зрения программной структуры серверная часть включает следующие подсистемы:\cite{fastapi_docs,sqlalchemy_docs,pydantic_docs}
\begin{enumerate}
  \item \texttt{auth} --- аутентификация, токены доступа и профиль пользователя;
  \item \texttt{tenants} --- работа с организациями и текущим контекстом организации-арендатора;
  \item \texttt{users} --- управление пользователями и их состоянием;
  \item \texttt{courses} и \texttt{lessons} --- управление курсами, уроками и прогрессом;
  \item \texttt{media} --- загрузка и выдача медиафайлов;
  \item \texttt{tests} --- жизненный цикл тестов и попыток;
  \item \texttt{adaptive} --- выбор следующего вопроса и расчёт рекомендаций;
  \item \texttt{recommendations} --- выдача списка рекомендаций слушателю;
  \item \texttt{analytics} --- агрегированная аналитика по системе.
\end{enumerate}

На рисунке~\ref{fig:design-components} приведена диаграмма компонентов серверной части. Она отражает декомпозицию серверного приложения на функциональные подсистемы и их связи с общими инфраструктурными механизмами.

\begin{figure}[H]
  \centering
  \includegraphics[width=\textwidth]{"UML component diagram"}
  \caption{Диаграмма компонентов системы}
  \label{fig:design-components}
\end{figure}

\section{Сценарии использования}

Во второй главе сценарии использования описываются как формальные последовательности действий, служащие основой для архитектурных и интерфейсных решений. Такое представление позволяет зафиксировать роли, предусловия, основной поток и результат выполнения для каждого ключевого сценария.\cite{grekul_proj_is,iso29119}

На основании диаграммы прецедентов и критического пути пользователей в таблице~\ref{tab:design-scenarios} приведено табличное представление ключевых сценариев, используемое для согласования ролей, предусловий и результатов выполнения.

\begin{table}[H]
\centering
\caption{Табличное представление основных сценариев использования}
\label{tab:design-scenarios}
{\setstretch{1}\fontsize{12pt}{12pt}\selectfont
\setlength{\tabcolsep}{4pt}
\setlength{\extrarowheight}{1pt}
\renewcommand{\arraystretch}{1}
\begin{tabularx}{\textwidth}{|>{\RaggedRight\arraybackslash}p{28mm}|>{\RaggedRight\arraybackslash}p{34mm}|>{\RaggedRight\arraybackslash}p{28mm}|>{\RaggedRight\arraybackslash}X|}
\hline
\textbf{Сценарий} & \textbf{Актор} & \textbf{Предусловие} & \textbf{Результат} \\
\hline
Подготовка обучения и назначение курса & Преподаватель, администратор организации & Пользователь аутентифицирован и работает в контексте своей организации & Курс создан, наполнен контентом и назначен слушателю. \\
\hline
Прохождение курса и адаптивного теста & Слушатель & Слушатель включён в организацию и имеет назначенный курс & Сформирован результат попытки и набор рекомендаций по слабым темам. \\
\hline
Анализ результатов обучения & Преподаватель, администратор организации & В системе накоплены данные по назначениям и попыткам & Получены агрегированные показатели и детализированные результаты по слушателям. \\
\hline
\end{tabularx}
}
\end{table}

\subsection{Сценарий 1: подготовка обучения и назначение курса}

\textbf{Акторы:} администратор организации, преподаватель.

\textbf{Предусловия:} пользователь аутентифицирован, организация выбрана, роль пользователя позволяет управлять контентом и назначениями.

\textbf{Основной поток:}
\begin{enumerate}
  \item преподаватель создаёт курс;
  \item преподаватель добавляет уроки и при необходимости медиа-контент;
  \item преподаватель создаёт тест и вопросы;
  \item администратор организации выбирает слушателя;
  \item администратор назначает курс слушателю;
  \item система фиксирует назначение и делает курс доступным в мобильном клиенте слушателя.
\end{enumerate}

\textbf{Постусловия:} слушатель получает назначенный курс и может приступить к изучению материала.

\subsection{Сценарий 2: прохождение курса и адаптивного теста}

\textbf{Актор:} слушатель.

\textbf{Предусловия:} слушатель аутентифицирован, имеет активное участие в организации, курс назначен, тест связан с курсом.

\textbf{Основной поток:}
\begin{enumerate}
  \item слушатель открывает список своих курсов;
  \item слушатель переходит в курс и изучает урок;
  \item слушатель запускает тест;
  \item система создаёт попытку с базовой сложностью;
  \item система выдаёт первый вопрос;
  \item слушатель отправляет ответ;
  \item система оценивает ответ и изменяет текущую сложность;
  \item система выдаёт следующий вопрос с учётом новой сложности;
  \item после достижения лимита вопросов система завершает попытку;
  \item система сохраняет результат, слабые темы и формирует рекомендации.
\end{enumerate}

\textbf{Постусловия:} результат попытки, история сложности и рекомендации доступны слушателю и административным ролям.

\subsection{Сценарий 3: анализ результатов обучения}

\textbf{Акторы:} преподаватель, администратор организации.

\textbf{Основной поток:}
\begin{enumerate}
  \item пользователь открывает раздел аналитики;
  \item система отображает агрегированные показатели по курсам, тестам и назначенным обучениям;
  \item пользователь открывает проблемные темы или индивидуальный отчёт по слушателю;
  \item система отображает результаты попыток, количество правильных ответов и рекомендации.
\end{enumerate}

\section{Модель данных}

Модель данных построена вокруг сущностей, необходимых для поддержки мультиарендного обучения (несколько организаций в одном экземпляре).\cite{date_intro_db_ru,stepanov_db,proprofs_multitenant_lms_2025,microsoft_multitenant_arch} На концептуальном уровне она отражает объекты предметной области и их семантические связи, а на логическом уровне фиксирует состав реляционных таблиц, ключей и ограничений целостности. Центральной сущностью является организация (\texttt{tenant}), относительно которой изолируются пользователи, курсы, уроки, тесты, попытки и рекомендации.

\subsection{Инфологическая модель базы данных}

Инфологическая модель описывает систему без привязки к конкретной СУБД и показывает укрупнённые сущности предметной области: организацию, пользователя, курс, урок, тему, тест, вопрос, попытку и результат. Такая степень абстракции позволяет наглядно показать логику движения данных от подготовки учебного контента до прохождения адаптивного теста и получения результата.

На рисунке~\ref{fig:design-infological} приведена инфологическая модель базы данных. На этом уровне намеренно опущены служебные и ассоциативные сущности, связанные с ролями, назначениями, рекомендациями и аудитом, так как они раскрываются в даталогической модели.

\begin{figure}[H]
  \centering
  \includegraphics[width=\textwidth]{db}
  \caption{Инфологическая модель базы данных}
  \label{fig:design-infological}
\end{figure}

\subsection{Даталогическая модель базы данных}

На даталогическом уровне инфологическая модель преобразуется в реляционную схему, реализованную средствами \texttt{SQLAlchemy}. При локальном запуске приложение может использовать \texttt{SQLite}, а при контейнеризированном развёртывании --- \texttt{PostgreSQL}; при этом набор таблиц, первичных и внешних ключей остаётся одинаковым. Для tenant-зависимых сущностей в схему включено поле \texttt{tenant\_id}, а для критичных связей заданы ограничения уникальности, исключающие дублирование членств, назначений и записей о прогрессе.

Для пояснительной записки даталогическая модель представлена двумя основными фрагментами, отражающими базовые контуры проектируемой системы: организационный контур и контур управления учебным контентом. Такой способ позволяет показать структуру таблиц и связей без перегрузки рисунков второстепенными служебными сущностями.

На рисунке~\ref{fig:design-datalogical-org} показана даталогическая модель организационного контура. Она описывает, каким образом в системе представлены организации, пользователи, роли, членства и учебные группы. Данный фрагмент показывает реализацию мультиарендного подхода: один пользователь может состоять в нескольких организациях, а его права определяются через таблицу членства и связанную с ней роль. Дополнительно здесь отражено объединение пользователей в группы, используемые при массовом назначении обучения.

\begin{figure}[H]
  \centering
  \includegraphics[width=\textwidth,height=0.42\textheight,keepaspectratio]{Даталогическая_организация}
  \caption{Даталогическая модель организационного контура}
  \label{fig:design-datalogical-org}
\end{figure}

На рисунке~\ref{fig:design-datalogical-courses} представлена даталогическая модель контура курсов. В неё включены таблицы, описывающие учебный контент, структуру курса, назначения, зачисления и фиксацию прогресса. Схема показывает, что курс принадлежит организации, состоит из уроков и тем, может назначаться как отдельному пользователю, так и группе, а факт прохождения отражается в таблицах зачислений и прогресса по урокам. Такой контур обеспечивает хранение как структуры учебных материалов, так и состояния их освоения слушателями.

\begin{figure}[H]
  \centering
  \includegraphics[width=\textwidth,height=0.56\textheight,keepaspectratio]{Даталогическая_курсы}
  \caption{Даталогическая модель контура курсов}
  \label{fig:design-datalogical-courses}
\end{figure}

Таким образом, в текущей версии пояснительной записки даталогическая модель раскрыта через два ключевых фрагмента: организационный и учебный. Первый описывает разграничение доступа и принадлежность пользователей к организациям и группам, а второй --- структуру курсов, механизм их назначения и фиксацию прогресса обучения. Контуры тестирования, рекомендаций и служебных таблиц при необходимости могут быть представлены отдельно без перегрузки основной части работы.

\section{Проектирование интерфейсов}

Проектирование интерфейсов выполнялось раздельно для административного и пользовательского контуров с учётом принципов проектирования взаимодействия.\cite{cooper2008aboutface,norman2013design,krug2014dontmake}

\subsection{Веб-интерфейс администратора и преподавателя}

Веб-клиент предназначен для работы с учебным контентом и аналитикой. Его основными разделами являются:
\begin{enumerate}
  \item вход в систему и выбор организации;
  \item управление пользователями;
  \item управление курсами;
  \item управление уроками и медиатекой;
  \item управление тестами;
  \item назначение курсов;
  \item просмотр аналитики;
  \item настройки профиля и языка интерфейса.
\end{enumerate}

Структура веб-интерфейса ориентирована на последовательное выполнение административных операций: сначала создаётся контент, затем выполняются назначения, после чего доступны данные аналитики.

\subsection{Мобильный интерфейс слушателя}

Мобильный клиент ориентирован на компактный пользовательский сценарий:
\begin{enumerate}
  \item вход пользователя;
  \item просмотр списка назначенных курсов;
  \item открытие курса и урока;
  \item запуск теста;
  \item просмотр результата, рекомендаций и истории попыток.
\end{enumerate}

На рисунках~\ref{fig:design-main-screen}--\ref{fig:design-courses-screen} показаны примеры экранов мобильного клиента.

\begin{figure}[H]
  \centering
  \setlength{\fboxrule}{0.4pt}
  \setlength{\fboxsep}{0pt}
  \fbox{\includegraphics[width=0.55\textwidth,height=0.75\textheight,keepaspectratio]{Главная}}
  \caption{Главный экран мобильного клиента}
  \label{fig:design-main-screen}
\end{figure}

\begin{figure}[H]
  \centering
  \setlength{\fboxrule}{0.4pt}
  \setlength{\fboxsep}{0pt}
  \fbox{\includegraphics[width=0.55\textwidth,height=0.75\textheight,keepaspectratio]{"Мои курсы - 1"}}
  \caption{Экран списка назначенных курсов}
  \label{fig:design-courses-screen}
\end{figure}

\section{Реализация web- и mobile-клиентов}

Клиентская часть системы реализована как два самостоятельных приложения, использующих единый REST API серверной части, но ориентированных на разные роли и сценарии работы. Веб-клиент предназначен для администратора и преподавателя, а мобильный клиент --- для слушателя. Такое разделение позволяет не перегружать интерфейсы лишними функциями и адаптировать взаимодействие с системой под реальные задачи пользователей.\cite{react_docs,flutter_docs}

На рисунке~\ref{fig:design-client-structure} представлена структурная схема клиентской части системы. Она показывает внутреннее устройство web- и mobile-приложений: пользовательский интерфейс, прикладные экраны, слой взаимодействия с API серверной части и локальное хранение пользовательского состояния. Общим для обоих клиентов является единый контракт обмена данными, а различие определяется набором экранов и способом управления состоянием.

\begin{figure}[H]
\centering
\resizebox{\textwidth}{!}{%
\begin{tikzpicture}[
  node distance=10mm and 20mm,
  every node/.style={font=\small},
  clientpanel/.style={
    draw,
    rounded corners,
    inner sep=4pt,
    align=left,
    fill=gray!5
  },
  apipanel/.style={
    draw,
    rounded corners,
    inner sep=4pt,
    align=left,
    fill=gray!10
  }
]
\node[clientpanel] (webclient) {%
  \begin{minipage}{54mm}\raggedright
  \textbf{Web-клиент}\par\smallskip\hrule\smallskip
  \texttt{React + TypeScript}, одностраничное приложение\par\smallskip\hrule\smallskip
  навигационный shell, административные страницы, формы CRUD и аналитика\par\smallskip\hrule\smallskip
  общий API-слой, токен доступа, код tenant, \texttt{localStorage}
  \end{minipage}
};
\node[clientpanel, below=of webclient] (mobileclient) {%
  \begin{minipage}{54mm}\raggedright
  \textbf{Mobile-клиент}\par\smallskip\hrule\smallskip
  \texttt{Flutter}, \texttt{MaterialApp} и \texttt{HomeShell}\par\smallskip\hrule\smallskip
  экраны курсов, уроков, адаптивного теста, рекомендаций и профиля\par\smallskip\hrule\smallskip
  \texttt{ApiClient}, состояние сессии в памяти, \texttt{SharedPreferences}
  \end{minipage}
};
\node[apipanel, right=20mm of $(webclient.east)!0.5!(mobileclient.east)$] (clientapi) {%
  \begin{minipage}{48mm}\raggedright
  \textbf{Единая точка интеграции}\par\smallskip\hrule\smallskip
  REST API, \texttt{Bearer}-авторизация, \texttt{X-Tenant-Code}\par\smallskip\hrule\smallskip
  общий контракт обмена данными для web- и mobile-клиента
  \end{minipage}
};

\coordinate (webclientout) at (webclient.east);
\coordinate (mobileclientout) at (mobileclient.east);
\coordinate (clientapiweb) at ($(clientapi.north west)!0.35!(clientapi.north east)$);
\coordinate (clientapimobile) at ($(clientapi.south west)!0.40!(clientapi.south east)$);
\coordinate (webclientbend) at (clientapiweb |- webclientout);
\coordinate (mobileclientbend) at (clientapimobile |- mobileclientout);
\draw[-{Latex[length=2.5mm]}, line width=0.8pt, rounded corners=2mm] ([xshift=1mm]webclientout) -- (webclientbend) -- ([yshift=0.8mm]clientapiweb);
\draw[-{Latex[length=2.5mm]}, line width=0.8pt, rounded corners=2mm] ([xshift=1mm]mobileclientout) -- (mobileclientbend) -- ([yshift=-0.8mm]clientapimobile);
\node[font=\footnotesize, fill=white, inner sep=1pt] at ($ ([xshift=1mm]webclientout)!0.5!(webclientbend) + (0, 3mm) $) {HTTPS/REST};
\node[font=\footnotesize, fill=white, inner sep=1pt] at ($ ([xshift=1mm]mobileclientout)!0.5!(mobileclientbend) + (0, -3mm) $) {HTTPS/REST};
\end{tikzpicture}
}
\caption{Структурная схема клиентской части системы}
\label{fig:design-client-structure}
\end{figure}

\subsection{Реализация web-клиента}

Веб-клиент реализован как одностраничное приложение на \texttt{React} и \texttt{TypeScript}. После аутентификации пользователь попадает в общий каркас приложения с боковой навигацией, из которого доступны основные разделы: панель показателей, переключение организации, пользователи, курсы, уроки, тесты, назначения, аналитика и настройки. Маршрутизация организована на стороне клиента, что позволяет быстро переходить между разделами без повторной загрузки страницы.

Работа с серверной частью вынесена в отдельный API-слой. Для всех запросов централизованно задаются базовый адрес сервера, заголовок авторизации \texttt{Bearer} и заголовок \texttt{X-Tenant-Code}, определяющий текущую организацию. Такой подход уменьшает дублирование кода и обеспечивает единое поведение для операций чтения, создания, редактирования и загрузки медиафайлов. Для типовых запросов используются отдельные функции отправки \texttt{GET}, \texttt{POST}, \texttt{PATCH} и \texttt{multipart/form-data}.

Сессионное состояние web-клиента включает токен доступа и код текущей организации. Эти данные передаются между страницами через состояние верхнего уровня приложения. Дополнительно в интерфейсе реализовано переключение языка, причём выбранный язык сохраняется в \texttt{localStorage} браузера и применяется при следующем открытии приложения.

Функциональные страницы web-клиента построены по единому принципу: загрузка данных через общий механизм запросов, отображение списка сущностей, форма создания или редактирования и визуальное подтверждение результата операции. Так реализованы управление пользователями, курсами, уроками, тестами и назначениями. В разделе уроков предусмотрено редактирование структуры страниц урока и загрузка медиафайлов, а в разделе аналитики --- загрузка агрегированных показателей, проблемных тем и данных по конкретному слушателю. В результате веб-клиент закрывает весь административный контур системы и напрямую поддерживает сценарии, описанные в первой главе.

\subsection{Реализация мобильного клиента}

Мобильный клиент реализован на \texttt{Flutter} как отдельное приложение слушателя. Корневой экран приложения построен на базе \texttt{MaterialApp}, а после входа пользователь переходит в основной \texttt{HomeShell}, содержащий нижнюю навигацию между тремя разделами: курсы, рекомендации и профиль. Для переключения между разделами используется \texttt{IndexedStack}, что позволяет сохранять состояние экранов при переходах.

Взаимодействие с сервером инкапсулировано в классе \texttt{ApiClient}. Он отвечает за формирование HTTP-запросов, добавление заголовков авторизации и заголовка организации (\texttt{X-Tenant-Code}), обработку ошибок и преобразование JSON-ответов в модели приложения. Через этот клиент выполняются аутентификация, загрузка курсов, получение структуры курса, получение рекомендаций, работа с попытками тестирования и загрузка данных для разбора завершённых попыток.

Состояние мобильного приложения разделено на два уровня. Языковые настройки сохраняются в \texttt{SharedPreferences}, благодаря чему выбор русского или английского языка восстанавливается после перезапуска приложения. Сессионные данные текущего пользователя передаются между экранами в памяти приложения, что достаточно для учебного минимально жизнеспособного продукта (MVP) и упрощает жизненный цикл клиента.

Экран списка курсов загружает назначенные пользователю курсы и открывает детальный экран курса. На экране курса выполняется параллельная загрузка структуры курса и связанных с ним тестов, после чего слушателю становятся доступны три действия: продолжение обучения, запуск адаптивного теста и просмотр истории попыток. Для прохождения урока используется отдельный экран-плеер, который получает данные урока с сервера и сохраняет прогресс пользователя, включая текущую страницу, набор завершённых страниц и позицию воспроизведения видео.

Отдельный экран \texttt{TestFlowScreen} реализует клиентскую часть адаптивного тестирования. Он инициирует попытку, получает очередной вопрос, отправляет выбранный ответ с учётом времени реакции и завершает попытку после исчерпания набора вопросов. После завершения теста слушатель видит результат, слабые темы и рекомендации. Дополнительно в мобильном клиенте реализованы экран рекомендаций и экран истории попыток с возможностью открыть детальный разбор завершённого теста. Таким образом, мобильный клиент не ограничивается просмотром контента, а поддерживает полный пользовательский цикл слушателя: от изучения курса до анализа собственных результатов.

\section{Реализация серверной логики и алгоритмов}

\subsection{Реализация API-подсистем}

Серверная часть реализует REST API, сгруппированное по назначению маршрутов. Такое разбиение определяет формат взаимодействия клиентов с серверной частью и набор доступных операций.
\begin{enumerate}
  \item \texttt{/auth} --- регистрация, вход, обновление токенов, получение профиля;
  \item \texttt{/tenants} --- список доступных организаций, текущая организация, выбор организации;
  \item \texttt{/users} --- просмотр, создание и изменение пользователей;
  \item \texttt{/courses} и \texttt{/lessons} --- работа с курсами, уроками, прогрессом и назначениями;
  \item \texttt{/media} --- библиотека и загрузка медиа;
  \item \texttt{/tests} и \texttt{/attempts} --- тесты, вопросы, попытки, история и разбор;
  \item \texttt{/recommendations} --- рекомендации слушателя;
  \item \texttt{/analytics} --- сводная и детальная аналитика.
\end{enumerate}

Такое разбиение соответствует структуре кода серверной части и обеспечивает независимое развитие функциональных подсистем при сохранении единой точки входа.

На рисунке~\ref{fig:design-server-structure} приведена структурная схема серверной части системы. В отличие от общей диаграммы компонентов, эта схема акцентирует именно уровни обработки запроса: маршрутный слой, контур tenant-контекста и разграничения доступа, сервисный слой бизнес-логики и слой хранения данных.

\begin{figure}[H]
\centering
\resizebox{\textwidth}{!}{%
\begin{tikzpicture}[
  every node/.style={font=\small},
  node distance=8mm,
  serverbox/.style={
    draw,
    rounded corners,
    rectangle split,
    rectangle split parts=3,
    rectangle split part align=left,
    text width=118mm,
    inner sep=4pt,
    fill=gray!5
  },
  storebox/.style={
    draw,
    rounded corners,
    text width=54mm,
    minimum height=11mm,
    align=center,
    fill=gray!10
  }
]
\node[serverbox] (routes) {\textbf{Маршрутный уровень}\nodepart{two}\texttt{/auth}, \texttt{/tenants}, \texttt{/users}, \texttt{/courses}, \texttt{/media}, \texttt{/tests}, \texttt{/recommendations}, \texttt{/analytics}\nodepart{three}принимает HTTP-запросы, валидирует входные данные и формирует ответы REST API};
\node[serverbox, below=of routes] (control) {\textbf{Контур доступа и tenant-контекста}\nodepart{two}tenant resolver, RBAC, проверка активного membership, аудит действий\nodepart{three}обеспечивает логическую изоляцию данных и контроль полномочий на каждом запросе};
\node[serverbox, below=of control] (services) {\textbf{Сервисный уровень}\nodepart{two}\texttt{auth}, \texttt{lesson\_player}, \texttt{adaptive}, \texttt{analytics}, \texttt{notifications}\nodepart{three}реализует бизнес-логику обучения, прохождения уроков, адаптивного тестирования и формирования рекомендаций};
\node[serverbox, below=of services] (data) {\textbf{Слой данных}\nodepart{two}SQLAlchemy-модели, схема БД, операции чтения и записи\nodepart{three}использует PK/FK, ограничения уникальности и \texttt{tenant\_id}-фильтрацию для tenant-зависимых сущностей};
\node[storebox, below left=10mm and 6mm of data] (dbstore) {\textbf{Реляционная БД}\\PostgreSQL / SQLite};
\node[storebox, below right=10mm and 6mm of data] (mediastore) {\textbf{Хранилище медиа}\\статические изображения и видео};

\draw[-{Latex[length=2.5mm]}] (routes) -- (control);
\draw[-{Latex[length=2.5mm]}] (control) -- (services);
\draw[-{Latex[length=2.5mm]}] (services) -- (data);
\draw[-{Latex[length=2.5mm]}] (data.south west) -- (dbstore.north);
\draw[-{Latex[length=2.5mm]}] (data.south east) -- (mediastore.north);
\end{tikzpicture}
}
\caption{Структурная схема серверной части системы}
\label{fig:design-server-structure}
\end{figure}

\subsection{Алгоритм адаптивного тестирования}

В проекте реализован дискретный алгоритм адаптивного тестирования, работающий в диапазоне сложностей от 1 до 5.\cite{lord1980,van_der_linden1997,wainer2000} Алгоритм основан на следующих правилах:
\begin{enumerate}
  \item при старте попытки текущая сложность равна базовой сложности теста;
  \item после правильного ответа сложность повышается на один уровень;
  \item после неправильного ответа сложность понижается на один уровень;
  \item следующее задание выбирается среди ещё не заданных вопросов, наиболее близких к текущей сложности;
  \item по ошибочным и медленным ответам накапливаются слабые темы;
  \item после завершения попытки формируются рекомендации по повторению материала.
\end{enumerate}

На рисунке~\ref{fig:design-adaptive-sequence} показана последовательность взаимодействия клиента и сервера при прохождении адаптивного теста. Диаграмма фиксирует обмен запросами между мобильным приложением, REST API серверной части и модулем адаптивной логики.

\begin{figure}[H]
\centering
\includegraphics[width=\textwidth,height=0.82\textheight,keepaspectratio]{sequence}
\caption{Последовательность прохождения адаптивного теста}
\label{fig:design-adaptive-sequence}
\end{figure}

Для иллюстрации логики изменения сложности на рисунке~\ref{fig:design-adaptive-flow} приведена схема алгоритма. Она отражает старт с базовой сложности, выбор ближайшего незаданного вопроса, изменение уровня сложности на шаг `+1/-1` и формирование рекомендаций после завершения попытки.

\begin{figure}[H]
\centering
\includegraphics[width=\textwidth,height=0.82\textheight,keepaspectratio]{"adaptive flow"}
\caption{Схема изменения сложности в адаптивном тестировании}
\label{fig:design-adaptive-flow}
\end{figure}

\subsection{Изоляция данных арендаторов}

Мультиарендная изоляция данных реализована на трёх уровнях:\cite{yojji_multitenant_lms_2025,owasp_top10_2021}
\begin{enumerate}
  \item на уровне запроса определяется код текущей организации;
  \item на уровне авторизации проверяется, что пользователь имеет активное участие в этой организации;
  \item на уровне доступа к данным выполняется фильтрация tenant-зависимых сущностей по полю \texttt{tenant\_id}.
\end{enumerate}

Таким образом, логическая изоляция не требует выделения отдельной базы данных на каждого арендатора и при этом обеспечивает корректное разделение данных в рамках минимально жизнеспособного продукта (MVP).

\section{Выводы по главе}

Во второй главе описана техническая реализация корпоративной LMS. Рассмотрены архитектура системы в виде модульного монолита, формальные сценарии использования, инфологическая и даталогическая модели базы данных (включая обоснование нормализации и связей между контурами организации, учебного контента и тестирования), проектирование интерфейсов веб-клиента и мобильного клиента, а также структурные схемы клиентской и серверной частей. Раскрыты состав REST API, алгоритмы адаптивного тестирования, формирования рекомендаций и логической изоляции данных арендаторов на уровне запроса, авторизации и доступа к данным. Процедуры контейнеризированного развёртывания, опытной эксплуатации и инструкции для пользователей и администраторов вынесены в третью главу и приложения пояснительной записки. Представленные решения соответствуют фактической реализации проекта и обеспечивают основу для демонстрации работоспособности системы на защите.\cite{grekul_proj_is,fowler_patterns_ru,booch2006uml}
